\documentclass[aspectratio=169,11pt]{beamer}

\usetheme{Madrid}
\usecolortheme{default}
\usefonttheme{professionalfonts}
\setbeamertemplate{navigation symbols}{}
\setbeamertemplate{footline}[frame number]

\usepackage{amsmath, amssymb, booktabs}

\newcommand{\E}{\mathbb{E}}
\newcommand{\1}{\mathbf{1}}
\newcommand{\Prb}{\mathbb{P}}

\title{Structural Modeling With Network Data}
\subtitle{Empirical Methods --- Week 20}
\author{Teaching deck draft}
\date{}

\begin{document}

\begin{frame}
  \titlepage
\end{frame}

\begin{frame}{Week 20 question}
\begin{itemize}
    \item When is a model of network formation or network-mediated behavior needed for counterfactuals?
    \item Lecture 18: define and diagnose the graph.
    \item Lecture 19: identify effects when links create exposure, interference, or dependence.
    \item Lecture 20: model policies that change links, information flows, search, matching, or equilibrium behavior.
\end{itemize}

\medskip
The central issue is not graph theory. It is whether the counterfactual changes the network.
\end{frame}

\begin{frame}{When reduced-form network methods are not enough}
\small
Reduced-form exposure design:
\[
Y_i = \alpha + \tau D_i(G,Z) + X_i'\beta + u_i
\]

Structural policy object:
\[
V(p)=\E\left[
W_i\left(G(p),a_i(G(p),X,\varepsilon),X_i\right)
\right]
\]

\begin{itemize}
    \item Reduced-form evidence can identify local exposure effects on the observed graph.
    \item It may not answer what happens when policy changes the graph itself.
    \item Examples: referral subsidies, platform ranking rules, team assignment, diffusion seeds, search access.
    \item The structural model earns its keep only if the counterfactual changes relationships or equilibrium behavior.
\end{itemize}
\end{frame}

\begin{frame}{Network methods block map}
\small
\begin{tabular}{p{0.24\linewidth} p{0.31\linewidth} p{0.33\linewidth}}
\toprule
Lecture & Main object & Failure if skipped \\
\midrule
18. Curation & nodes, edges, weights, timing, boundaries & unclear estimand and fragile measurement \\
19. Causal inference & exposure, interference, dependence & peer correlations mistaken for effects \\
20. Structure & formation, behavior, counterfactuals & link-changing policies overclaimed \\
\bottomrule
\end{tabular}

\medskip
The three layers are complements. The structural model should inherit the graph discipline from Lecture 18 and the design discipline from Lecture 19.
\end{frame}

\begin{frame}{Network formation vs behavior}
\small
\begin{columns}[T,onlytextwidth]
\begin{column}{0.48\linewidth}
\textbf{Behavior on a network}
\[
a_i
=
\alpha+\beta x_i
+\gamma\sum_j g_{ij}a_j
+\delta\sum_j g_{ij}x_j+\varepsilon_i
\]
\begin{itemize}
    \item Conditions on observed links.
    \item Needs timing, equilibrium, and exposure logic.
    \item Useful for peer effects, diffusion, learning, and adoption.
\end{itemize}
\end{column}
\begin{column}{0.48\linewidth}
\textbf{Network formation}
\[
G_{ij}
=
\1\{U_{ij}(X_i,X_j,G_{-ij},\eta_{ij})\ge 0\}
\]
\begin{itemize}
    \item Explains why links exist.
    \item Allows policy to change links.
    \item Needs assumptions on opportunity, sorting, and strategic payoffs.
\end{itemize}
\end{column}
\end{columns}
\end{frame}

\begin{frame}{Structural peer effects and equilibrium}
\small
\[
a_i
=
\alpha
+ \beta x_i
+ \gamma \sum_j g_{ij}a_j
+ \delta \sum_j g_{ij}x_j
+ \varepsilon_i
\]

\begin{itemize}
    \item The equation is familiar; the structural burden is the equilibrium interpretation.
    \item Are peer actions simultaneous or lagged?
    \item Is the graph fixed before actions, or chosen jointly with actions?
    \item Can multiple equilibria matter?
    \item Which variation separates peer effects from sorting and common shocks?
\end{itemize}

\medskip
A coefficient is not yet a counterfactual model.
\end{frame}

\begin{frame}{Search, referrals, and labor-market applications}
\small
\[
\lambda_i(G)=\lambda_0+\lambda_1\sum_j G_{ij}q_j
\]

\begin{itemize}
    \item Links can carry vacancy information, screening signals, trust, reputation, or favoritism.
    \item Referral policies may change who forms links, who receives information, and how employers update beliefs.
    \item Matching policies can change feasible worker--firm pairs and congestion.
    \item Inequality can rise or fall depending on whether new links bridge opportunity gaps or reinforce segregation.
\end{itemize}

\medskip
Labor anchors: Calvo-Armengol and Jackson; Dustmann et al.; Galenianos; Pallais and Sands.
\end{frame}

\begin{frame}{Diffusion and equilibrium behavior}
\small
\[
\Pr(a_{it}=1)
=
\sigma\left(
\alpha+X_i'\beta+\rho\sum_j g_{ij}a_{j,t-1}+\kappa s_{it}
\right)
\]

\begin{itemize}
    \item Diffusion models study how information, technologies, practices, and beliefs spread.
    \item Relevant for training, management practices, job-search norms, platform information, and wage knowledge.
    \item Equilibrium response matters when adoption changes congestion, incentives, or search behavior.
    \item Validation requires adoption paths, not only final adoption rates.
\end{itemize}
\end{frame}

\begin{frame}{Estimation, computation, and validation}
\small
\[
m(\theta)=\E_\theta[
\text{degree},
\text{clustering},
\text{homophily},
\text{wage covariance across links},
\text{job transitions},
\text{referral yields}]
\]

\begin{itemize}
    \item Use moments or likelihood objects that discipline the counterfactual primitives.
    \item Formation moments: degree, homophily, clustering, link persistence.
    \item Behavior moments: outcomes conditional on links, transitions, adoption paths, referral yields.
    \item Computation often uses simulation, local statistics, sparse approximations, or separable blocks.
    \item Validation should compare fixed-network and endogenous-link counterfactuals.
\end{itemize}
\end{frame}

\begin{frame}{Policy counterfactual discipline}
\small
\begin{tabular}{p{0.26\linewidth} p{0.31\linewidth} p{0.31\linewidth}}
\toprule
Policy changes & Need to model & Validation target \\
\midrule
Existing-link information & behavior on fixed graph & outcome response by exposure \\
New referral links & formation and referral yields & degree, homophily, callback rates \\
Platform ranking & matching and congestion & search sets, applications, matches \\
Training seeds & diffusion and equilibrium response & adoption timing and spillovers \\
Team assignment & peer groups and sorting & group composition and productivity \\
\bottomrule
\end{tabular}

\medskip
The policy defines the necessary structure.
\end{frame}

\begin{frame}{Network Methods block summary}
\small
\begin{itemize}
    \item \textbf{Where the literature is now:} richer administrative and platform networks; sharper exposure and interference designs; more structural formation, search, referral, and diffusion models.
    \item \textbf{Key open questions:} endogenous links, dynamic networks, privacy-preserving network data, reduced-form discipline for structural counterfactuals, and distributional effects of network policies.
    \item \textbf{Research opportunities:} referral hiring, internal labor markets, coworker learning, manager networks, platform matching, school-to-work contacts, and local opportunity networks.
\end{itemize}

\medskip
The frontier is not bigger graphs. It is better alignment between mechanism, data, design, model, and counterfactual.
\end{frame}

\begin{frame}{Research Lab design}
\begin{columns}[T,onlytextwidth]
\begin{column}{0.32\linewidth}
\textbf{Reproduce}
\begin{itemize}
    \item synthetic workers
    \item dyadic opportunities
    \item link formation model
    \item network moments
    \item behavior equation
\end{itemize}
\end{column}
\begin{column}{0.32\linewidth}
\textbf{Diagnose}
\begin{itemize}
    \item topology fit
    \item behavior fit
    \item moment relevance
    \item fixed vs endogenous
    \item validation memo
\end{itemize}
\end{column}
\begin{column}{0.32\linewidth}
\textbf{Transfer}
\begin{itemize}
    \item referral search
    \item cross-group links
    \item match access
    \item inequality
    \item policy memo
\end{itemize}
\end{column}
\end{columns}
\end{frame}

\begin{frame}{Takeaways}
\begin{itemize}
    \item Structural network models are useful when policy changes links, information flows, matching sets, or equilibrium behavior.
    \item Behavior-on-network models and formation models answer different questions.
    \item Endogenous links make assumptions and validation burdens heavier.
    \item Moments must be chosen for the counterfactual, not for convenience.
    \item Reduced-form network evidence remains valuable as discipline for structural claims.
    \item The Network Methods block gives a sequence: define, identify, then model.
\end{itemize}
\end{frame}

\end{document}
